\newglossaryentry{Wirtschaftliche_Bevoelkerung}
{   name={Wirtschaftliche Wohnbevölkerung},
    description={Die Wohnbevölkerung umfast alle Personen, die in der Stadt Zürich gemeldet sind,         in der Stadt wohnen und die städtische Infrastruktur beanspruchen. Die Wohnbevölkerung            umfast zudem Wochenendaufenthalterinnen und Wochenendaufenthalter, Asylsuchende,                  Flüchtlinge mit vorläufiger Aufnahme sowie kurzfristige Aufenhtalterinnen und Aufenthalter         . Das Bevölkerungsregister der Stadt Zürich beinhaltet die wirtschaftliche Bevölkerung}
}
\newglossaryentry{Konversionssaldo}{
    name={Konversionssaldo},
    description={Zahl der Konversionen zu einer Religion abzüglich der Konversionen weg von dieser         Religion}
}
\newacronym{BFS}{BFS}{Bundesamt für Statistik}
\newacronym{ABS}{ABS}{Alternative Bank Schweiz}
\newglossaryentry{Konfidenzintervall}{
    name={Konfidenzintervall},
    description={Das 95-Prozent-Konfidenzintervall bezeichnet den Bereich, der bei unendlicher                     Wiederholung eines Zufallsexperiments mit einer Wahrscheinlichkeit von 95 Prozent                  den wahren Wert der Grundgesamtheit einschliesst. Das Konfidenzintervall wird auch 
    Vertrauensintervall oder Erwartungsbereich genannt}
}
\newglossaryentry{KonfidenzintervallA}{
    name={Konfidenzintervall A},
    description={Das 95-Prozent-Konfidenzintervall bezeichnet den Bereich, der bei unendlicher                     Wiederholung eines Zufallsexperiments mit einer Wahrscheinlichkeit von 95 Prozent                  den wahren Wert der Grundgesamtheit einschliesst. Das Konfidenzintervall wird auch 
    Vertrauensintervall oder Erwartungsbereich genannt}
}
\newglossaryentry{KonfidenzintervallB}{
    name={Konfidenzintervall B},
    description={Das 95-Prozent-Konfidenzintervall bezeichnet den Bereich, der bei unendlicher                     Wiederholung eines Zufallsexperiments mit einer Wahrscheinlichkeit von 95 Prozent                  den wahren Wert der Grundgesamtheit einschliesst. Das Konfidenzintervall wird auch 
    Vertrauensintervall oder Erwartungsbereich genannt}
}
